\section{第12章\quad 曲面积分}
\subsection{第一类曲面积分}
\begin{exercise}{520}{基础·上半球对称性}
设 $S$ 为单位上半球面。利用对称性与球面参数计算 $\iint_S(x+y+z)\,\dd S$。
\end{exercise}
\begin{exercise}{521}{基础·平面面积元素}
设 $S:z=2y-1,x^2+(y-1)^2\leq1$。写出面积元素并计算 $\iint_Sz\,\dd S$。
\end{exercise}
\begin{exercise}{522}{基础·平面约束化简}
设 $S$ 为 $x/2+y/3+z/4=1$ 在第一卦限的部分。利用平面方程计算 $\iint_S(2x+4y/3+z)\,\dd S$。
\end{exercise}
\begin{exercise}{523}{基础·平面三角形积分}
设 $S=\{x+y+z=1,x,y,z\geq0\}$。计算 $\iint_Sy^2\,\dd S$。
\end{exercise}
\begin{exercise}{524}{基础·球面约束}
设 $S$ 为单位上半球面。利用球面方程计算 $\iint_Sz(x^2+y^2+z^2+1)\,\dd S$。
\end{exercise}
\begin{exercise}{525}{基础·球冠积分}
设 $S=\{x^2+y^2+z^2=1,z\geq1/2\}$。用球面参数求 $\iint_Sz\,\dd S$。
\end{exercise}
\begin{exercise}{526}{常规·分片曲面面积}
设 $S:|x|+|y|+|z|=1$。利用八个全等三角形求曲面面积及 $\iint_Sx\,\dd S$。
\end{exercise}
\begin{exercise}{527}{基础·球面二阶矩}
设 $S:x^2+y^2+z^2=R^2$，$R>0$。利用坐标互换求 $\iint_Sz^2\,\dd S$。
\end{exercise}
\begin{exercise}{528}{基础·锥面对称性}
设 $S:z=\sqrt{x^2+y^2},0\leq z\leq1$。利用对称性求 $\iint_S(xy+yz)\,\dd S$。
\end{exercise}
\begin{exercise}{529}{基础·锥面积分}
设 $S:z=\sqrt{x^2+y^2},0\leq z\leq1$。写出面积元素并计算 $\iint_Sz\,\dd S$。
\end{exercise}
\begin{exercise}{530}{基础·半球平均高度}
设 $S$ 为半径 $a>0$、球心在原点的上半球面。求 $\iint_Sz\,\dd S$ 及均匀曲面的平均高度。
\end{exercise}
\begin{exercise}{531}{常规·锥体边界分片}
设 $S$ 为 $\{\sqrt{x^2+y^2}\leq z\leq1\}$ 的完整边界。分成锥面与顶圆盘，求 $\iint_S(x^2+y^2)\,\dd S$。
\end{exercise}
\begin{exercise}{532}{基础·平面曲面多项式积分}
设 $S=\{x+y+z=1,x,y,z\geq0\}$。计算 $\iint_Sxy\,\dd S$。
\end{exercise}
\begin{exercise}{533}{常规·投影与一元换元}
设 $S=\{x+y+z=1,x,y,z\geq0\}$。计算 $\iint_S(1+x+y)^{-2}\,\dd S$。提示：投影后令 $u=x+y,v=x$。
\end{exercise}
\begin{exercise}{534}{基础·圆柱侧面积分}
设 $S:x^2+y^2=4,0\leq z\leq3$。计算 $\iint_S(x^2+y^2)^{-1}\,\dd S$。
\end{exercise}
\begin{exercise}{535}{基础·锥面参数化}
设 $S:z=\sqrt3\sqrt{x^2+y^2},0\leq z\leq\sqrt3$。写出参数和面积元素，求 $\iint_S(x^2+y^2)\,\dd S$。
\end{exercise}
\begin{exercise}{536}{基础·抛物面面积}
设 $S:z=1-x^2-y^2,z\geq0$。计算曲面面积。
\end{exercise}
\begin{exercise}{537}{常规·抛物面面积积分}
设 $S:z=x^2+y^2,x^2+y^2\leq1$。用极坐标计算 $\iint_Sz\,\dd S$。
\end{exercise}
\begin{exercise}{538}{基础·平面投影积分}
设 $S$ 为平面 $2x+2y+z=6$ 在第一卦限的部分。写出投影区域及面积元素，求 $\iint_Sx\,\dd S$。
\end{exercise}
\begin{exercise}{539}{基础·球冠对称积分}
设 $S=\{x^2+y^2+z^2=4,z\geq1\}$。计算 $\iint_S(x+y+z)\,\dd S$。
\end{exercise}
\begin{exercise}{540}{基础·球面积分变式}
设 $S$ 为单位上半球面。利用球面参数计算 $\iint_Sz^2\,\dd S$。
\end{exercise}
\begin{exercise}{541}{常规·锥面质量与重心}
设 $S:z=\sqrt{x^2+y^2},0\leq z\leq1$，面密度为 $\rho=z$。求质量与重心高度。
\end{exercise}
\subsection{第二类曲面积分}
\begin{exercise}{542}{基础·球面内侧投影}
设 $S$ 为单位球面在第一卦限的部分，取内侧。将 $\iint_Sz\,\dd x\dd y$ 化为平面二重积分并计算。
\end{exercise}
\begin{exercise}{543}{基础·平面下侧}
设 $S:z=0,x^2+y^2\leq1$，取下侧。计算 $\iint_S(x^2+y^2)\,\dd x\dd y$。
\end{exercise}
\begin{exercise}{544}{基础·下半球投影}
设 $S:x^2+y^2+z^2=1,z\leq0$，取下侧。计算 $\iint_Sx^2\,\dd x\dd y$，并说明改取上侧时的结果。
\end{exercise}
\begin{exercise}{545}{基础·上半球通量投影}
设 $S$ 为单位上半球面，取上侧。计算 $\iint_Sz^2\,\dd x\dd y$。
\end{exercise}
\begin{exercise}{546}{基础·闭柱体分面通量}
设 $S$ 为 $x^2+y^2\leq1,1\leq z\leq2$ 的完整边界，取外侧。分别求侧面及上下底面的 $\iint_Sz\,\dd x\dd y$，再求总和。
\end{exercise}
\begin{exercise}{547}{基础·四分之一柱面定向}
设 $S:x^2+y^2=1,x,y\geq0,0\leq z\leq3$，取外侧。计算 $\iint_Sx\,\dd y\dd z+y\,\dd z\dd x+z\,\dd x\dd y$。
\end{exercise}
\begin{exercise}{548}{基础·两类曲面积分转换}
设 $S=\{x+y+z=1,x,y,z\geq0\}$，取上侧。求单位法向量，将连续场 $F=(P,Q,R)$ 的通量写成第一类曲面积分。
\end{exercise}
\begin{exercise}{549}{基础·图像曲面通量}
设 $S:z=1-x^2-y^2,z\geq0$，取上侧。将 $F=(2x,2y,1)$ 的通量化为圆盘积分并计算。
\end{exercise}
\begin{exercise}{550}{基础·常法向量}
设 $S=\{x-y+z=1,x\geq0,y\leq0,z\geq0\}$，取上侧。求 $F=(1,2,1)$ 通过 $S$ 的通量。
\end{exercise}
\begin{exercise}{551}{基础·闭柱面无奇点通量}
设 $S$ 为 $x^2+y^2\leq1,-1\leq z\leq1$ 的完整边界，取外侧。分面计算 $\iint_Sx\,\dd y\dd z+z^2\,\dd x\dd y$。
\end{exercise}
\begin{exercise}{552}{基础·柱面无效分量}
设 $S:x^2+y^2=4,0\leq z\leq2$，取外侧。计算 $\iint_Sx\,\dd y\dd z+y\,\dd z\dd x+(z+1)\,\dd x\dd y$。
\end{exercise}
\begin{exercise}{553}{基础·抛物面下侧}
设 $S:z=(x^2+y^2)/2,0\leq z\leq2$，取下侧。计算 $\iint_Sz\,\dd x\dd y$。
\end{exercise}
\begin{exercise}{554}{常规·半球法向量}
设 $S$ 为单位上半球面，取外侧。利用球面单位法向量求 $F=(xz,yz,z^2)$ 的通量。
\end{exercise}
\begin{exercise}{555}{基础·柱面参数积分}
设 $S:x^2+y^2=1,0\leq z\leq2$，取外侧。用参数计算 $\iint_Sxz\,\dd y\dd z$。
\end{exercise}
\begin{exercise}{556}{基础·抛物面方向}
设 $S:z=x^2+y^2,z\leq1$，取下侧。求 $\iint_Sx^2\,\dd y\dd z+z\,\dd x\dd y$。
\end{exercise}
\begin{exercise}{557}{基础·四分之一柱面}
设 $S:x^2+y^2=1,x,y\geq0,0\leq z\leq1$，取外侧。求 $\iint_Sy\,\dd y\dd z$。
\end{exercise}
\begin{exercise}{558}{基础·斜面下侧}
设 $S=\{x+y-z=1,x,y\geq0,z\leq0\}$，取下侧。写出投影区域，计算 $\iint_S\dd y\dd z+z\,\dd x\dd y$。
\end{exercise}
\begin{exercise}{559}{基础·无散场通量}
设 $S$ 为单位圆柱体 $x^2+y^2\leq1,0\leq z\leq1$ 的边界外侧。计算场 $F=(y-z,z-x,x-y)$ 的总通量。
\end{exercise}
\begin{exercise}{560}{基础·球面径向通量}
设 $S$ 为半径 $R>0$、球心在原点的球面外侧。用单位法向量求场 $F=(x,y,z)$ 的通量。
\end{exercise}
\begin{exercise}{561}{常规·四面体分面通量}
设 $S$ 为四面体 $x,y,z\geq0,x+y+z\leq1$ 的边界外侧。求 $F=(yz,zx,xy)$ 的总通量，再由三个坐标面通量求斜面通量。
\end{exercise}
\begin{exercise}{562}{常规·球面径向奇点场}
设 $S:x^2+y^2+z^2=4$，取外侧。直接用球面法向量计算 $\iint_S\frac{x\,\dd y\dd z+y\,\dd z\dd x+z\,\dd x\dd y}{(x^2+y^2+z^2)^{3/2}}$，说明在球体内部直接用高斯公式缺少的条件。
\end{exercise}
\begin{exercise}{563}{基础·下半球奇偶性}
设 $S$ 为单位下半球面，取下侧。计算 $\iint_Sx^2\,\dd y\dd z+y^2\,\dd z\dd x+z^2\,\dd x\dd y$。
\end{exercise}
\begin{exercise}{564}{基础·补面计算柱面通量}
设 $S:x^2+y^2=1,0\leq z\leq2$，取外侧。补上两个底面，用高斯公式求 $F=(x,y,z)$ 的侧面通量。
\end{exercise}
\begin{exercise}{565}{常规·锥面通量}
设 $S:z=1-\sqrt{x^2+y^2},0\leq z\leq1$，取上侧。补上底面，求 $F=(x,y,2z)$ 的通量。
\end{exercise}
\begin{exercise}{566}{基础·立方体通量}
设 $S=\partial[0,1]^3$，取外侧。用高斯公式求 $F=(xy,y^2,z)$ 的总通量。
\end{exercise}
\begin{exercise}{567}{基础·切向场变式}
设 $S:z=\sqrt{x^2+y^2},1\leq z\leq2$，取下侧。写出一个法向量，说明 $F=(x,y,z)$ 通过 $S$ 的通量为零。
\end{exercise}
\begin{exercise}{568}{常规·抛物面通量小综合}
设 $S:z=1-x^2-y^2,z\geq0$，取上侧。求 $F=(x,y,z)$ 的通量，并用补面高斯公式核验。
\end{exercise}
\subsection{高斯公式}
\begin{exercise}{569}{基础·球体外通量}
设 $S$ 为单位球面外侧。用高斯公式求 $\iint_Sx\,\dd y\dd z+y\,\dd z\dd x+z\,\dd x\dd y$，写明使用公式的条件。
\end{exercise}
\begin{exercise}{570}{常规·锥面补面}
设 $S:z=\sqrt{x^2+y^2},0\leq z\leq1$，取下侧。补上顶圆盘，计算 $F=(x,2y,3(z-1))$ 通过 $S$ 的通量。
\end{exercise}
\begin{exercise}{571}{基础·上半球补面}
设 $S$ 为单位上半球面外侧。用补面法求 $F=(x,y,z+1)$ 的通量，明确底面方向。
\end{exercise}
\begin{exercise}{572}{基础·常散度通量}
设 $S$ 为 $0\leq z\leq1-x^2-y^2$ 的完整边界外侧。求场 $F=(y,z,x+z)$ 的总通量。
\end{exercise}
\begin{exercise}{573}{基础·半球体二次矩}
设 $S$ 为单位上半球体的完整边界外侧。用高斯公式求 $F=(x^3,y^3,z^3)$ 的总通量。
\end{exercise}
\begin{exercise}{574}{基础·全边界与侧面}
设圆柱体 $\Omega=\{x^2+y^2\leq1,0\leq z\leq3\}$。分别求场 $F=(0,0,z)$ 通过完整边界和侧面的外通量。
\end{exercise}
\begin{exercise}{575}{基础·立方体与对称性}
分别求 $F=(x^2,y^2,z^2)$ 通过 $[0,1]^3$ 与 $[-1,1]^3$ 边界外侧的通量。
\end{exercise}
\begin{exercise}{576}{基础·球面多项式通量}
设 $S$ 为半径 $2$、球心在原点的球面外侧。用高斯公式求 $F=(x^3,y^3,z^3)$ 的总通量。
\end{exercise}
\begin{exercise}{577}{基础·通量恢复体积}
设 $S$ 是围成区域 $\Omega$ 的分片光滑闭曲面，取外侧，已知 $F=(x,y,z)$ 的通量为 $12$。求 $\Omega$ 的体积。
\end{exercise}
\begin{exercise}{578}{基础·单位立方体}
设 $S=\partial[0,1]^3$，取外侧。求 $F=(4xz,-y^2,yz)$ 的通量。
\end{exercise}
\begin{exercise}{579}{常规·斜顶柱体}
设 $\Omega=\{x^2+y^2\leq1,0\leq z\leq1-x\}$。用高斯公式求场 $F=(x,y,z)$ 通过完整边界的外通量。
\end{exercise}
\begin{exercise}{580}{基础·绕轴旋转面}
设 $S:y=x^2+z^2,0\leq y\leq1$，取法向量的 $y$ 分量为负的一侧。补上 $y=1$ 的圆盘，求场 $F=(x,0,z)$ 的通量。
\end{exercise}
\begin{exercise}{581}{基础·非零补面通量}
设 $S$ 为半径 $2$、球心在原点的上半球面外侧。用补面法求场 $F=(0,yz,2)$ 的通量。
\end{exercise}
\begin{exercise}{582}{基础·方向余弦通量}
设 $S$ 为单位球面外侧，$n=(\cos\alpha,\cos\beta,\cos\gamma)$。用高斯公式求 $\iint_S(x^3\cos\alpha+y^3\cos\beta+z^3\cos\gamma)\,\dd S$。
\end{exercise}
\begin{exercise}{583}{基础·锥面方向余弦}
设 $S:z=\sqrt{x^2+y^2},0\leq z\leq1$，取下侧。写出单位法向量并计算 $F=(0,0,z^2)$ 的通量。
\end{exercise}
\begin{exercise}{584}{常规·开曲面常数项}
设 $S:z=1-x^2-y^2,z\geq0$，取上侧。用补面法求 $F=(x,y,z-1)$ 的通量。
\end{exercise}
\begin{exercise}{585}{常规·椭圆投影}
设 $S:z=1-x^2-y^2/4,z\geq0$，取上侧。求场 $F=(x,y,0)$ 的通量，明确椭圆投影换元的雅可比因子。
\end{exercise}
\begin{exercise}{586}{基础·对称项消去}
设 $S:z=x^2+y^2,z\leq1$，取上侧。用直接投影法求 $F=(x^2,y^2,z-1)$ 的通量。
\end{exercise}
\begin{exercise}{587}{基础·四面体通量}
设 $S$ 为 $x,y,z\geq0,2x+y+2z\leq2$ 的完整边界外侧。用高斯公式求 $F=(x^2+1,-2y,3z)$ 的通量。
\end{exercise}
\begin{exercise}{588}{常规·底面非零的补面法}
设 $S:z=1-x^2-y^2,z\geq0$，取上侧。用补面法求 $F=(-x,-y,x^2-z)$ 的通量。
\end{exercise}
\begin{exercise}{589}{基础·抛物体闭边界}
设 $S$ 为 $x^2+y^2\leq z\leq1$ 的完整边界外侧。用高斯公式求 $F=(x,y,z^2)$ 的通量。
\end{exercise}
\begin{exercise}{590}{基础·四面体坐标矩}
设 $S$ 为 $x,y,z\geq0,x+y+z\leq1$ 的边界外侧。求 $\iint_Sxy\,\dd y\dd z+yz\,\dd z\dd x+xz\,\dd x\dd y$。
\end{exercise}
\begin{exercise}{591}{常规·球锥区域通量}
设 $\Omega=\{x^2+y^2+z^2\leq1,z\geq\sqrt{x^2+y^2}\}$。用球面坐标求场 $F=(x,y,z)$ 通过完整边界的外通量。
\end{exercise}
\begin{exercise}{592}{基础·锥顶连续场}
设 $S:z=\sqrt{x^2+y^2},0\leq z\leq1$，取下侧。用补面法求 $F=(x,y,e^z)$ 的通量。
\end{exercise}
\begin{exercise}{593}{基础·球壳通量}
设 $S$ 为球壳 $1\leq x^2+y^2+z^2\leq4$ 的完整边界，取相对于球壳的外侧。求 $F=(x,y,z)$ 的总通量。
\end{exercise}
\begin{exercise}{594}{基础·球壳方向变式}
设 $S$ 为 $1\leq\sqrt{x^2+y^2+z^2}\leq2$ 的完整边界，取外侧。分别求常向量场 $F=(0,0,1)$ 通过内外球面的通量，并求总和。
\end{exercise}
\begin{exercise}{595}{常规·高斯公式小综合}
设 $S$ 为单位球面外侧，$F=(ax,y,z)$，$a\in\mathbb R$。求通量，并求使通量为零的 $a$；说明场的光滑性与积分面闭合性。
\end{exercise}
\subsection{斯托克斯公式}
\begin{exercise}{596}{基础·三角形旋度积分}
设 $L$ 为 $x+y+z=1$ 在第一卦限的三角形边界，与上侧法向符合右手定则。用斯托克斯公式求 $\oint_Ly\,\dd x+z\,\dd y+x\,\dd z$。
\end{exercise}
\begin{exercise}{597}{基础·竖直圆盘选面}
设 $L:y^2+z^2=1,x=0$，方向与法向量 $(1,0,0)$ 符合右手定则。计算 $\oint_L-y\,\dd z+z\,\dd y$。
\end{exercise}
\begin{exercise}{598}{基础·三角形边界方向}
设 $L$ 沿 $(1,0,0)\to(0,1,0)\to(0,0,1)\to(1,0,0)$。用斯托克斯公式计算 $\oint_Lz\,\dd x+x\,\dd y+y\,\dd z$，写出所取法向量。
\end{exercise}
\begin{exercise}{599}{基础·水平正方形环量}
设 $L$ 为 $0\leq x,y\leq1,z=1$ 的边界，从正 $z$ 轴向下看为逆时针。求 $\oint_L(y-z)\,\dd x+(z-x)\,\dd y$。
\end{exercise}
\begin{exercise}{600}{常规·球面大圆环量}
设 $L:x^2+y^2+z^2=1,x+y+z=0$，方向与法向量 $(1,1,1)$ 符合右手定则。用平面圆盘求 $\oint_Ly\,\dd x+z\,\dd y+x\,\dd z$。
\end{exercise}
\begin{exercise}{601}{基础·水平截圆}
设 $L:x^2+y^2=4,z=2$，从正 $z$ 轴向下看为逆时针。求 $\oint_L3y\,\dd x-xz\,\dd y+yz^2\,\dd z$。
\end{exercise}
\begin{exercise}{602}{基础·精确项与闭路}
设 $L:x^2+y^2=9,z=0$ 取从正 $z$ 轴向下看为逆时针的方向。求 $\oint_L2y\,\dd x+3x\,\dd y-z^2\,\dd z$。
\end{exercise}
\begin{exercise}{603}{常规·同边界换面}
设 $A=(y^2,xy,xz)$，$S$ 为单位上半球面上侧。用斯托克斯公式将旋度通量转到边界或单位圆盘，求 $\iint_S(\nabla\times A)\cdot n\,\dd S$，写明圆盘方向。
\end{exercise}
\begin{exercise}{604}{常规·五面曲面方向}
设 $A=(-y,x,0)$，$S$ 为单位立方体边界外侧去掉 $z=0$ 底面后的五个面。求 $\iint_S(\nabla\times A)\cdot n\,\dd S$，说明底部边界从正 $z$ 轴向下看的方向。
\end{exercise}
\begin{exercise}{605}{基础·常旋度三角形}
设 $L$ 沿 $(1,0,0)\to(0,1,0)\to(0,0,1)\to(1,0,0)$。求 $\oint_L(z-y)\,\dd x+(x-z)\,\dd y+(y-x)\,\dd z$。
\end{exercise}
\begin{exercise}{606}{基础·斜圆盘投影}
设 $L:x^2+y^2=1,z=1-x$，从正 $z$ 轴向下看为逆时针。用斯托克斯公式求 $\oint_L-y\,\dd x+x\,\dd y+\dd z$。
\end{exercise}
\begin{exercise}{607}{基础·空间参数环量}
设 $L:x=\cos t,y=\sin t,z=\cos t$，$t:0\to2\pi$。求 $\oint_Ly\,\dd x+z\,\dd y+x\,\dd z$。
\end{exercise}
\begin{exercise}{608}{基础·反向边界功}
设 $F=(y,z,x)$，$L$ 是 $x+y+z=1$ 第一卦限三角形的边界，方向与下侧法向符合右手定则。用斯托克斯公式求功，并给出反向后的值。
\end{exercise}
\begin{exercise}{609}{基础·同边界旋度通量}
设 $A\in C^1(\mathbb R^3)$，$S_1$ 为单位上半球面上侧，$S_2$ 为单位圆盘上侧。利用斯托克斯公式说明两曲面上的旋度通量相等。
\end{exercise}
\begin{exercise}{610}{常规·截圆环量小综合}
设 $L:x^2+y^2+z^2=4,z=1$，从正 $z$ 轴向下看为逆时针。求场 $A=(-yz,xz,0)$ 的环量，并用参数法核验符号。
\end{exercise}
\subsection{曲面积分的应用}
\begin{exercise}{611}{基础·无散场总通量}
设 $A=(yz,xz,xy)$，$S$ 为单位圆柱体 $x^2+y^2\leq1,0\leq z\leq1$ 的完整边界外侧。求散度与总通量。
\end{exercise}
\begin{exercise}{612}{基础·多项式通量}
设 $S=\partial[0,1]^3$，取外侧。求场 $A=(2x-z,x^2y,-xz^2)$ 的通量。
\end{exercise}
\begin{exercise}{613}{基础·平移球面常散度}
设 $S$ 为球心 $(3,-1,2)$、半径 $3$ 的球面外侧。求 $A=(2x+3z,-xz-y,y^2+2z)$ 的通量。
\end{exercise}
\begin{exercise}{614}{常规·半圆柱通量}
设 $S:y^2+z^2=1,z\geq0,0\leq x\leq1$，取柱体外侧。求 $A=(0,yz,z^2)$ 的通量。
\end{exercise}
\begin{exercise}{615}{基础·六面通量分配}
求 $A=(x,y,z)$ 通过单位立方体 $[0,1]^3$ 六个面的各自外通量与总通量。
\end{exercise}
\begin{exercise}{616}{基础·散度与球面对称性}
求 $A=(x^2+yz,y^2+xz,z^2+xy)$ 的散度，并利用对称性求单位球面外通量。
\end{exercise}
\begin{exercise}{617}{基础·复合函数散度}
求 $A=(e^{xy},\cos(xy),\cos(xz^2))$ 的散度及其在 $(0,1,0)$ 处的值。
\end{exercise}
\begin{exercise}{618}{基础·散度与旋度}
对 $A=(y^2,xy,xz)$ 分别求 $\nabla\cdot A$ 与 $\nabla\times A$。
\end{exercise}
\begin{exercise}{619}{基础·半圆柱体总通量}
设 $\Omega=\{0\leq x\leq1,y^2+z^2\leq1,z\geq0\}$。用高斯公式求 $A=(0,yz,z^2)$ 通过完整边界的外通量。
\end{exercise}
\begin{exercise}{620}{基础·点散度与局部通量}
设 $A=(xy^2,ye^z,x\ln(1+z^2))$，$S_r$ 为球心 $(1,1,0)$、半径 $r$ 的球面外侧。求 $\lim_{r\to0^+}\frac{\iint_{S_r}A\cdot n\,\dd S}{4\pi r^3/3}$。提示：先求散度再用积分中值定理。
\end{exercise}
\begin{exercise}{621}{基础·旋度与圆周环量}
求 $A=(x^2-y,4z,x^2)$ 的旋度，再求其沿 $z=0$ 平面内逆时针单位圆周的环量。
\end{exercise}
\begin{exercise}{622}{基础·梯度场势函数}
设 $A=(z+\sin y,x\cos y-z,x-y)$。验证其旋度为零，并求一个势函数。
\end{exercise}
\begin{exercise}{623}{基础·三角函数场旋度}
求 $A=(x^2\sin y,y^2\sin(xz),xy\sin z)$ 的旋度及其在原点的值。
\end{exercise}
\begin{exercise}{624}{基础·平面圆周环量}
设 $A=(xy,-yz,xz)$，$L$ 为 $z=0$ 平面上以 $(1,1,0)$ 为圆心、半径 $1$ 的圆周，从正 $z$ 轴向下看为逆时针。用斯托克斯公式求环量。
\end{exercise}
\begin{exercise}{625}{基础·径向函数梯度}
在 $r=\sqrt{x^2+y^2+z^2}>0$ 上令 $u=\ln r$。求 $\nabla u$ 与 $\nabla\cdot\nabla u$。
\end{exercise}
\begin{exercise}{626}{常规·径向函数拉普拉斯}
设 $u=1/\sqrt{x^2+y^2+z^2}$，定义于原点以外。直接求 $\nabla u$，并检验 $\nabla\cdot\nabla u=0$。
\end{exercise}
\begin{exercise}{627}{基础·梯度场无旋}
设 $u$ 有连续二阶偏导数。利用混合偏导数相等逐分量证明 $\nabla\times\nabla u=0$。
\end{exercise}
\begin{exercise}{628}{基础·散度旋度恒等式}
设 $A=(P,Q,R)$ 有连续二阶偏导数。逐项展开并证明 $\nabla\cdot(\nabla\times A)=0$。
\end{exercise}
\begin{exercise}{629}{常规·梯度与球面通量}
设 $u=x^2+y^2+z^2$。求 $\nabla u$，再用高斯公式求其通过半径 $2$、球心在原点的球面外侧的通量。
\end{exercise}
\subsubsection{常规综合练习}
\begin{exercise}{630}{基础·球面法向量}
设 $S:x^2+y^2+z^2=4$。求外单位法向量，并计算 $\iint_Sz^2\,\dd S$。
\end{exercise}
\begin{exercise}{631}{基础·柱面积分}
设 $S:x^2+y^2=1,0\leq z\leq1$。计算 $\iint_S(1+z^2)^{-1}\,\dd S$。
\end{exercise}
\begin{exercise}{632}{基础·球面平方函数}
设 $S$ 为单位球面。利用对称性计算 $\iint_S(x+y+1)^2\,\dd S$。
\end{exercise}
\begin{exercise}{633}{基础·平移球面积分}
设 $S:x^2+y^2+z^2=2z$。求球心、半径，再计算 $\iint_Sz\,\dd S$。
\end{exercise}
\begin{exercise}{634}{基础·球冠面积}
求单位球面位于 $z\geq\sqrt3/2$ 部分的面积。
\end{exercise}
\begin{exercise}{635}{常规·分片球面函数}
设 $S$ 为单位球面，$f=z$（$z\geq1/2$），其余处 $f=0$。计算 $\iint_Sf\,\dd S$。
\end{exercise}
\begin{exercise}{636}{基础·下半球取上侧}
设 $S$ 为单位下半球面，取上侧。计算 $\iint_S(1+z)\,\dd x\dd y$，说明该方向与球体外侧的关系。
\end{exercise}
\begin{exercise}{637}{常规·球面奇点场}
设 $S$ 为单位球面外侧，$r=\sqrt{x^2+y^2+z^2}$。直接计算 $F=(x,y,z)/r^3$ 的通量，说明原点奇点为何妨碍在整个球体上直接应用高斯公式。
\end{exercise}
\begin{exercise}{638}{基础·椭圆曲面投影}
设 $S:z=1-x^2/4-y^2,z\geq0$，取上侧。计算 $\iint_Sz\,\dd x\dd y$，写出椭圆换元。
\end{exercise}
\begin{exercise}{639}{常规·补面方向核验}
设 $S:z=x^2+y^2,z\leq1$，取上侧。用高斯公式求 $F=(x,y,z-1)$ 的通量，说明 $S$ 作为立体 $x^2+y^2\leq z\leq1$ 边界时的外侧方向。
\end{exercise}
\begin{exercise}{640}{基础·零散度参数}
设 $F=(ax,y,-3z)$。求使 $F$ 的散度恒为零的 $a$，并说明此时它通过任意分片光滑闭曲面外侧的通量。
\end{exercise}
\begin{exercise}{641}{常规·锥台侧面}
设 $S:z=\sqrt{x^2+y^2},1\leq z\leq2$，取相对于锥体的侧面外侧。补上两个圆盘，求 $F=(y,-x,z^2)$ 的通量。
\end{exercise}
\begin{exercise}{642}{基础·斜截柱面环量}
设 $L:x^2+y^2=1,x-y+z=2$，从正 $z$ 轴向下看为逆时针。用斯托克斯公式求 $\oint_L-y\,\dd x+x\,\dd y$。
\end{exercise}
\begin{exercise}{643}{基础·水平圆周环量}
设 $L:z=x^2+y^2,z=4$，从正 $z$ 轴向下看为逆时针。计算 $\oint_L2y\,\dd x+xz\,\dd y-yz^2\,\dd z$。
\end{exercise}
\begin{exercise}{644}{基础·平面约束化简}
设 $L:x^2+y^2=1,y+z=0$，从正 $z$ 轴向下看为逆时针。用 $z=-y$ 计算 $\oint_Lz\,\dd x+y\,\dd z$。
\end{exercise}
\begin{exercise}{645}{常规·斜圆盘旋度积分}
设 $L:x^2+y^2=1,z=x+y$，从正 $z$ 轴向下看为逆时针。用斯托克斯公式求 $\oint_Lxz\,\dd x+x\,\dd y+y^2\,\dd z/2$。
\end{exercise}
\begin{exercise}{646}{基础·菱形空间边界}
设 $L:x+y+z=2,|x|+|y|=1$，从正 $z$ 轴向下看为逆时针。求 $\oint_L-y\,\dd x+x\,\dd y+\dd z$。
\end{exercise}
\begin{exercise}{647}{常规·开曲线补闭合}
设 $L:x=\cos t,y=\sin t,z=0$，$t:0\to\pi$。补上直径线段，用斯托克斯公式求 $\int_L(y+z)\,\dd x+x\,\dd y+\dd z$。
\end{exercise}
\begin{exercise}{648}{基础·半球面前侧}
设 $S:x=\sqrt{1-y^2-z^2}$，取法向量 $x$ 分量非负的一侧。用投影法求 $F=(1,0,0)$ 的通量。
\end{exercise}
\begin{exercise}{649}{常规·抛物面通量应用}
设 $S:z=1+x^2+y^2,1\leq z\leq2$，取上侧。求 $F=(x,y,z)$ 的通量，写出投影区域与定向面积元素。
\end{exercise}
\begin{exercise}{650}{基础·平移球面通量}
设 $S:(x-2)^2+y^2+z^2=1$，取外侧。用高斯公式求 $F=(x,y,z)$ 的通量。
\end{exercise}
\begin{exercise}{651}{基础·曲面约束分母}
设 $S:y=x^2+z^2,0\leq y\leq1$，取法向量 $y$ 分量为负的一侧。先用曲面方程化简分母，再求 $F=(0,1,0)/(y-x^2-z^2+1)$ 的通量。
\end{exercise}
\begin{exercise}{652}{常规·旋转曲面通量}
将线段 $x=0,y=1,0\leq z\leq1$ 绕 $z$ 轴旋转得到侧面 $S$，取外侧。写出曲面方程并求 $F=(x,y,z^2)$ 的通量。
\end{exercise}
\begin{exercise}{653}{基础·锥面切向场变式}
设 $S:z=\sqrt{x^2+y^2},1\leq z\leq2$，取下侧。利用法向量求 $F=(y,-x,0)$ 的通量。
\end{exercise}
\begin{exercise}{654}{常规·斯托克斯小综合}
设 $S$ 为平面 $x+y+z=1$ 在第一卦限的部分，取上侧，其边界 $L$ 按右手定则定向。用斯托克斯公式求 $\oint_L(y-\cos z)\,\dd x+z\,\dd y+(x+x\sin z)\,\dd z$；提示：含三角函数的两项组成 $\dd(-x\cos z)$。
\end{exercise}
